% !TeX program = pdflatex
\documentclass[10pt, aspectratio=169]{beamer}

\input{../../_en_preambulo_beamer}
% Comandos y macros estadísticas compartidas para las presentaciones Beamer en inglés (Metropolis)

% Conjuntos numéricos básicos
\newcommand{\R}{\mathbb{R}}
\newcommand{\N}{\mathbb{N}}
\newcommand{\Q}{\mathbb{Q}}
\newcommand{\Z}{\mathbb{Z}}
\newcommand{\set}[1]{\left\{ #1 \right\}}
\newcommand{\abs}[1]{\left|#1\right|}
\newcommand{\norm}[1]{\left\| #1 \right\|}

% Comandos estadísticos y probabilísticos del curso
\newcommand{\Var}{\operatorname{Var}}
\newcommand{\cov}{\operatorname{Cov}}
\newcommand{\corr}{\rho}
\newcommand{\comb}[2]{\begin{pmatrix} #1 \\ #2 \end{pmatrix}}
\newcommand{\s}{\sigma}
\newcommand{\particion}{\mathfrak{P}}
\newcommand{\card}[1]{\#\left( #1 \right)}
\newcommand{\seq}[3]{\set{#1}_{#2}^{#3}}
\newcommand{\E}{\mathbb{E}}
\newcommand{\Prob}{\mathbb{P}}

% Entornos pedagógicos y cajas en inglés académico natural (soportando tanto nombres en inglés como en español)
\newenvironment{discussion}{%
  \begin{alertblock}{Discussion Question}%
}{%
  \end{alertblock}%
}
\newenvironment{discusion}{%
  \begin{alertblock}{Discussion Question}%
}{%
  \end{alertblock}%
}

\newenvironment{reflection}{%
  \begin{alertblock}{Classroom Reflection}%
}{%
  \end{alertblock}%
}
\newenvironment{reflexion}{%
  \begin{alertblock}{Classroom Reflection}%
}{%
  \end{alertblock}%
}

\newenvironment{paradox}{%
  \begin{alertblock}{Exploratory Paradox}%
}{%
  \end{alertblock}%
}
\newenvironment{paradoja}{%
  \begin{alertblock}{Exploratory Paradox}%
}{%
  \end{alertblock}%
}

\newenvironment{motivation}{%
  \begin{exampleblock}{Motivation: Data Science in Practice}%
}{%
  \end{exampleblock}%
}
\newenvironment{motivacion}{%
  \begin{exampleblock}{Motivation: Data Science in Practice}%
}{%
  \end{exampleblock}%
}

\newenvironment{quickchallenge}{%
  \begin{block}{Quick Team Challenge}%
}{%
  \end{block}%
}
\newenvironment{retoclase}{%
  \begin{block}{Quick Team Challenge}%
}{%
  \end{block}%
}


\title[Discrete Probability Mass Functions]{Section 03.01: Discrete Probability Mass Functions}
\subtitle{From Abstract Sample Spaces to the Real Line: Support and PMF Axioms}
\author[J. Castillo Colmenares]{Juliho Castillo Colmenares}
\institute{Tecnológico de Monterrey}
\date{\vspace{-2.0cm}}

\begin{document}

% Slide 01: Cover
\begin{frame}[plain]
  \vspace{-0.5cm}
  \titlepage
\end{frame}

% Slide 02: Roadmap
\begin{frame}{Roadmap — Chapter 03: Discrete Random Variables}
  \begin{block}{Curricular Structure of Chapter 03}
    \footnotesize
    \begin{enumerate}\itemsep=0.08cm
      \item \textbf{\textcolor{TecRojo}{Section 03.01: Discrete Probability Mass Functions (PMF and Support)}} $\leftarrow$ \emph{Current focus}
      \item \textcolor{gray}{Section 03.02: Cumulative Distribution Functions (Discrete CDF)}
      \item \textcolor{gray}{Section 03.03: Mathematical Expectation, Variance, and Moments}
      \item \textcolor{gray}{Section 03.04: Bernoulli and Binomial Distributions}
      \item \textcolor{gray}{Section 03.05: Geometric and Negative Binomial Distributions}
      \item \textcolor{gray}{Section 03.06: Hypergeometric Distribution}
    \end{enumerate}
  \end{block}
  \vspace{-0.05cm}
  \pause
  \begin{alertblock}{Learning Objective of Section 03.01}
    \footnotesize
    Formalize the mapping from abstract sample spaces $\Omega$ to scalar probability distributions on $\mathbb{R}$, characterizing the measurable support $\mathcal{S}_X$, mastering Kolmogorov's axioms for discrete variables ($\sum p_X(x)=1$), and performing non-linear variable transformations.
  \end{alertblock}
\end{frame}

% Slide 03: Motivation — From Sample Space to the Real Line
\begin{frame}{Motivation — From Abstract Sample Space to the Real Line $\mathbb{R}$}
  \begin{columns}[T]
    \begin{column}{0.48\textwidth}
      \begin{block}{Limitations of Abstract Categories $\Omega$}
        \small
        \begin{itemize}\itemsep=0.1cm
          \item Real-world random phenomena usually yield qualitative or symbolic outcomes:
          \begin{itemize}
            \item Coin toss sequences ($\{HHT, THH, \dots\}$).
            \item Diagnostic status ($\{\text{Healthy}, \text{Failure}\}$).
            \item Network packet headers or server logs.
          \end{itemize}
          \item Differential calculus and linear algebra cannot operate directly on words or symbolic strings.
        \end{itemize}
      \end{block}
    \end{column}
    \begin{column}{0.48\textwidth}
      \begin{alertblock}{The Solution: Quantitative Projection}
        \small
        \begin{itemize}\itemsep=0.1cm
          \item We define a measurable function that assigns a unique real number $X(\omega) \in \mathbb{R}$ to every empirical outcome $\omega \in \Omega$.
          \item This projection extracts the exact operational metric of interest (error counts, response times, item yields).
          \item It converts set theory into high-precision functional and numerical analysis.
        \end{itemize}
      \end{alertblock}
    \end{column}
  \end{columns}
\end{frame}

% Slide 04: Formal Definition of Discrete Random Variable and Support
\begin{frame}{Formal Definition of Discrete Random Variable and Support}
  \fontsize{7.5pt}{9pt}\selectfont
  \begin{block}{Discrete Random Variable (D.R.V.)}
    \fontsize{7.5pt}{9pt}\selectfont
    A random variable $X: \Omega \to \mathbb{R}$ is called \textbf{discrete} if the set of all its possible values is finite or countably infinite. Formally, its range can be put into one-to-one correspondence with a subset of the natural numbers $\mathbb{N}$.
  \end{block}
  \vspace{-0.15cm}
  \pause
  \begin{alertblock}{The Measurable Support $\mathcal{S}_X$}
    \fontsize{7.5pt}{9pt}\selectfont
    The \textbf{support} of a discrete random variable $X$, denoted by $\mathcal{S}_X$, is the precise set of real numbers where probability mass is concentrated with strictly positive weight:
    $$\mathcal{S}_X = \{x \in \mathbb{R} \mid P(X = x) > 0\} = \{x_1, x_2, x_3, \dots\}$$
    \begin{itemize}\itemsep=0.08cm
      \item Outside the support $\mathcal{S}_X$, the probability is identically zero: $P(X = x) = 0 \quad \forall x \notin \mathcal{S}_X$.
      \item Any subset $A \subset \mathbb{R}$ is evaluated solely by examining its intersection with $\mathcal{S}_X$.
    \end{itemize}
  \end{alertblock}
\end{frame}

% Slide 05: Probability Mass Function (PMF)
\begin{frame}{Probability Mass Function (PMF)}
  \fontsize{7.5pt}{9pt}\selectfont
  \begin{block}{Analytical Characterization of Point Mass}
    \fontsize{7.5pt}{9pt}\selectfont
    For a discrete random variable $X$, its \textbf{Probability Mass Function} (PMF), denoted by $p_X(x)$ or $f_X(x)$, is the real-valued function $p_X: \mathbb{R} \to [0, 1]$ defined by:
    $$p_X(x) = P(X = x) = P(\{\omega \in \Omega \mid X(\omega) = x\})$$
    This relation associates each real number with the probability of the preimage of all elementary events in $\Omega$ that collapse onto the value $x$.
  \end{block}
  \vspace{-0.15cm}
  \pause
  \begin{alertblock}{Indicator Function Representation}
    \fontsize{7.5pt}{9pt}\selectfont
    To prevent ambiguities in continuous domains or analytical extensions, the PMF is algebraically expressed using the support indicator function $\mathbf{1}_{\mathcal{S}_X}(x)$:
    $$p_X(x) = g(x) \cdot \mathbf{1}_{\mathcal{S}_X}(x) = \begin{cases}
      g(x) & \text{if } x \in \mathcal{S}_X \\
      0 & \text{if } x \notin \mathcal{S}_X
    \end{cases}$$
  \end{alertblock}
\end{frame}

% Slide 06: The Two Kolmogorov Axioms in Discrete Probability
\begin{frame}{The Two Kolmogorov Axioms in Discrete Probability}
  \begin{block}{Structural Axioms of Admissibility}
    A real-valued function $p_X(x)$ constitutes a legitimate probability mass function for a discrete random variable if and only if it simultaneously satisfies two fundamental measure-theoretic axioms:
    \begin{enumerate}\itemsep=0.2cm
      \item \textbf{Strict Non-Negativity Axiom:} Probability cannot assume negative values at any point on the real line:
      $$p_X(x) \ge 0 \quad \text{for all } x \in \mathbb{R}$$
      \item \textbf{Unitary Summation Axiom (Normalization):} The sum of all probability weights distributed across the countable support must equal exactly one:
      $$\sum_{x_k \in \mathcal{S}_X} p_X(x_k) = 1$$
    \end{enumerate}
  \end{block}
\end{frame}

% Slide 07: Graphical Anatomy and Event Evaluation Theorem
\begin{frame}{Graphical Anatomy and Event Evaluation Theorem}
  \begin{columns}[T]
    \begin{column}{0.48\textwidth}
      \begin{block}{Graphical Anatomy — Mass Spectrum}
        \small
        \begin{itemize}\itemsep=0.1cm
          \item Graphically, a PMF is depicted as a vertical bar chart or point impulses (discrete Dirac delta functions).
          \item The height of each bar at point $x_k$ along the horizontal axis measures the exact scalar value of $p_X(x_k)$.
          \item No probability mass exists in the intermediate gaps between support points.
        \end{itemize}
      \end{block}
    \end{column}
    \begin{column}{0.48\textwidth}
      \begin{alertblock}{Universal Evaluation Theorem}
        \small
        \begin{itemize}\itemsep=0.1cm
          \item For any measurable event or subset $A \subset \mathbb{R}$, the probability that $X$ falls within $A$ is obtained by direct summation:
          $$P(X \in A) = \sum_{x_k \in A \cap \mathcal{S}_X} p_X(x_k)$$
          \item This allows evaluating intervals of arbitrary topology by summing interior support points.
        \end{itemize}
      \end{alertblock}
    \end{column}
  \end{columns}
\end{frame}

% Slide 08: Non-Linear Variable Transformations ($Y=g(X)$)
\begin{frame}{Non-Linear Variable Transformations ($Y=g(X)$)}
  \fontsize{7.5pt}{9pt}\selectfont
  \begin{block}{The Operational Mapping Problem}
    \fontsize{7.5pt}{9pt}\selectfont
    Often, we are not directly interested in the original variable $X$, but rather in a transformed quantity, such as power dissipation ($Y = X^2$), absolute deviation ($Y = |X - \mu|$), or logarithmic scaling ($Y = \ln X$). Let $g: \mathbb{R} \to \mathbb{R}$ be any arbitrary mapping and define $Y = g(X)$.
  \end{block}
  \vspace{-0.15cm}
  \pause
  \begin{alertblock}{Preimage Theorem and Induced Support}
    \fontsize{7.5pt}{9pt}\selectfont
    The transformed random variable $Y$ is also discrete. Its induced support $\mathcal{S}_Y$ is the direct image of the original support: $\mathcal{S}_Y = \{y = g(x) \mid x \in \mathcal{S}_X\}$. To obtain the induced PMF $p_Y(y)$, we sum the probabilities of all preimages that map to each specific $y$:
    $$p_Y(y) = P(Y = y) = \sum_{x \in \mathcal{S}_X: g(x) = y} p_X(x)$$
    If $g(x)$ is non-injective (e.g., absolute values or even powers), multiple distinct points of $\mathcal{S}_X$ aggregate into a single point in $\mathcal{S}_Y$.
  \end{alertblock}
\end{frame}

% Slide 09A: Python Lab — Normalization and Polynomial Support
\begin{frame}[fragile]{Python Lab — Normalization and Polynomial Support (1/4)}
  \begin{block}{Source Code: \texttt{verify\_polynomial\_pmf} (\texttt{03.01\_pmf\_and\_support.py})}
    \fontsize{6pt}{7pt}\selectfont
    \lstinputlisting[language=Python, firstline=10, lastline=28]{../../code/03_variables_aleatorias_discretas/03.01_pmf_and_support.py}
  \end{block}
\end{frame}

% Slide 09B: Python Lab — Monte Carlo Simulation (Dice)
\begin{frame}[fragile]{Python Lab — Monte Carlo Simulation of Dice Sums (2/4)}
  \begin{block}{Source Code: \texttt{simulate\_dice\_sum\_pmf} (\texttt{03.01\_pmf\_and\_support.py})}
    \fontsize{6pt}{7pt}\selectfont
    \lstinputlisting[language=Python, firstline=30, lastline=45]{../../code/03_variables_aleatorias_discretas/03.01_pmf_and_support.py}
  \end{block}
\end{frame}

% Slide 09C: Python Lab — Non-Linear Transformation
\begin{frame}[fragile]{Python Lab — Non-Linear Support Transformation (3/4)}
  \begin{block}{Source Code: \texttt{verify\_transformed\_pmf} (\texttt{03.01\_pmf\_and\_support.py})}
    \fontsize{6pt}{7pt}\selectfont
    \lstinputlisting[language=Python, firstline=47, lastline=61]{../../code/03_variables_aleatorias_discretas/03.01_pmf_and_support.py}
  \end{block}
\end{frame}

% Slide 09D: Python Lab — Verified Console Output
\begin{frame}[fragile]{Python Lab — Verified Console Output (4/4)}
  \begin{block}{Autoverified Terminal Execution (\texttt{Terminal Output})}
    \fontsize{6.5pt}{7.5pt}\selectfont
\begin{verbatim}
=== 1. Polynomial PMF Verification ===
Support S_X: [1 2 3 4]
Normalization Constant c: 0.033333 (Exact: 1/30)
Total Probability Sum: 1.000000
P(X is even | X >= 2): 0.6897 (Exact: 20/29)

=== 2. Monte Carlo Dice Sum Validation ===
Trials: 100,000
Max Empirical vs Exact PMF Error: 0.00168
Empirical P(Sum=7): 0.1669 (Exact: 0.1667)

=== 3. Transformation Y = |X - 1| ===
Induced Support S_Y: [0 1 2 3]
  P(Y = 0) = 0.30
  P(Y = 1) = 0.30
  P(Y = 2) = 0.25
  P(Y = 3) = 0.15
Sum of Induced PMF: 1.00
\end{verbatim}
  \end{block}
\end{frame}

% Slide 10A: Classroom Exercise — Tier 1: Fundamental (Statement)
\begin{frame}{Classroom Exercise — Tier 1: Fundamental (1/2)}
  \begin{block}{Problem 3.1.3 — Kolmogorov Axiomatic Audit}
    \footnotesize
    Consider a candidate function $g: \mathbb{Z} \to \mathbb{R}$ defined over the set of integers by the algebraic expression:
    $$g(x) = \begin{cases}
      \dfrac{x^2 - 1}{10} & \text{if } x \in \{-2, -1, 0, 1, 2\}, \\[8pt]
      0 & \text{otherwise.}
    \end{cases}$$
    \begin{enumerate}\itemsep=0.15cm
      \item By evaluating point by point across the set $\{-2, -1, 0, 1, 2\}$, verify whether the function satisfies the \textbf{Strict Non-Negativity Axiom}.
      \item Calculate the global sum over $\mathbb{Z}$ to determine whether it satisfies the \textbf{Unitary Normalization Axiom}. Conclude regarding its admissibility as a PMF.
    \end{enumerate}
  \end{block}
\end{frame}

% Slide 10B: Classroom Exercise — Tier 1: Fundamental (Solution)
\begin{frame}{Classroom Exercise — Tier 1: Fundamental (2/2)}
  \fontsize{6.8pt}{8.2pt}\selectfont
  \begin{block}{1. Audit of the Non-Negativity Axiom ($g(x) \ge 0$)}
    \fontsize{6.8pt}{8.2pt}\selectfont
    We directly evaluate each of the $5$ points where the formula takes a non-trivial value:
    \begin{itemize}\itemsep=0.06cm
      \item $x = -2 \implies g(-2) = \frac{(-2)^2 - 1}{10} = \frac{3}{10} = 0.30 \ge 0$.
      \item $x = -1 \implies g(-1) = \frac{(-1)^2 - 1}{10} = \frac{0}{10} = 0.00 \ge 0$.
      \item $x = 0 \implies g(0) = \frac{0^2 - 1}{10} = -\frac{1}{10} = \mathbf{-0.10 < 0}$.
      \item $x = 1 \implies g(1) = \frac{1^2 - 1}{10} = 0.00 \ge 0$.
      \item $x = 2 \implies g(2) = \frac{2^2 - 1}{10} = \frac{3}{10} = 0.30 \ge 0$.
    \end{itemize}
    By yielding a strictly negative value at $x=0$, $g(x)$ \textbf{blatantly violates the first axiom} ($P(X=0) \ge 0$).
  \end{block}
  \vspace{-0.15cm}
  \pause
  \begin{alertblock}{2. Audit of the Normalization Axiom ($\sum g(x) = 1$)}
    \fontsize{6.8pt}{8.2pt}\selectfont
    Summing across all candidate support points: $\displaystyle \sum_{x=-2}^{2} g(x) = \frac{3}{10} + 0 - \frac{1}{10} + 0 + \frac{3}{10} = \frac{5}{10} = 0.50 \neq 1$.
    \textbf{Conclusion:} The function $g(x)$ is inadmissible as a PMF because it violates both Kolmogorov axioms simultaneously.
  \end{alertblock}
\end{frame}

% Slide 11A: Classroom Exercise — Tier 2: Operational (Statement)
\begin{frame}{Classroom Exercise — Tier 2: Operational (1/2)}
  \begin{block}{Problem 3.1.6 — Maximum of Two Drawn Spheres}
    \footnotesize
    An urn contains $6$ identical spheres numbered $1$ through $6$. Two spheres are drawn successively and without replacement. Let $X_1$ be the number on the first sphere and $X_2$ the number on the second. We define the discrete random variable $M = \max(X_1, X_2)$ as the maximum of the two numbers obtained.
    \begin{enumerate}\itemsep=0.15cm
      \item Derive analytically the probability mass function $f_M(m) = P(M = m)$ for each $m \in \{2, 3, 4, 5, 6\}$ by analyzing preimages in the sample space.
      \item Calculate the exact probability that the observed maximum is greater than or equal to $5$ ($P(M \ge 5)$).
    \end{enumerate}
  \end{block}
\end{frame}

% Slide 11B: Classroom Exercise — Tier 2: Operational (Solution)
\begin{frame}{Classroom Exercise — Tier 2: Operational (2/2)}
  \begin{columns}[T]
    \begin{column}{0.48\textwidth}
      \begin{block}{1. Derivation of the Maximum PMF}
        \fontsize{6.5pt}{7.8pt}\selectfont
        The sample space $\Omega$ contains $|\Omega| = 6 \times 5 = 30$ equiprobable ordered pairs $(x_1, x_2)$ with $x_1 \neq x_2$. For the maximum to equal precisely $m$:
        \begin{itemize}\itemsep=0.06cm
          \item The larger sphere is $m$ and the other is strictly smaller ($1 \le x < m$).
          \item There are $m-1$ choices for the smaller sphere, occurring in $2$ permutations $((m,x)$ or $(x,m))$ $\implies 2(m-1)$ ordered pairs out of $30$:
          $$\displaystyle f_M(m) = \frac{2(m-1)}{30} = \frac{m-1}{15}, \quad m \in \{2 \dots 6\}$$
        \end{itemize}
      \end{block}
    \end{column}
    \begin{column}{0.48\textwidth}
      \begin{alertblock}{2. Table and Upper Tail Calculation}
        \fontsize{6.5pt}{7.8pt}\selectfont
        Evaluating the PMF table over $\mathcal{S}_M = \{2, 3, 4, 5, 6\}$:
        \begin{itemize}\itemsep=0.06cm
          \item $f_M(2) = 1/15$, $f_M(3) = 2/15$, $f_M(4) = 3/15$.
          \item $f_M(5) = 4/15$, $f_M(6) = 5/15$.
          \item Total sum: $(1+2+3+4+5)/15 = 15/15 = 1$.
        \end{itemize}
        Calculating the upper tail probability $P(M \ge 5)$:
        $$\displaystyle P(M \ge 5) = f_M(5) + f_M(6) = \frac{4}{15} + \frac{5}{15} = \mathbf{\frac{9}{15} = \frac{3}{5} = 0.60}$$
      \end{alertblock}
    \end{column}
  \end{columns}
\end{frame}

% Slide 12A: Classroom Exercise — Tier 3: Analytical (Statement)
\begin{frame}{Classroom Exercise — Tier 3: Analytical (1/2)}
  \begin{block}{Problem 3.1.7 — Normalization via Telescoping Sums}
    \footnotesize
    Let $N \in \mathbb{N}$ be a fixed positive integer. Consider the family of discrete random variables $X_N$ with finite support $\mathcal{S}_N = \{1, 2, \dots, N\}$ whose probability mass function takes the rational form:
    $$f_N(x) = \frac{c_N}{x(x+1)}, \quad x \in \{1, 2, \dots, N\}$$
    where $c_N$ is a real normalization constant to be determined.
    \begin{enumerate}\itemsep=0.15cm
      \item Prove rigorously, using partial fraction decomposition and telescoping summation, that the required normalization constant is exactly:
      $$c_N = \frac{N+1}{N}$$
    \end{enumerate}
  \end{block}
\end{frame}

% Slide 12B: Classroom Exercise — Tier 3: Analytical (Solution)
\begin{frame}{Classroom Exercise — Tier 3: Analytical (2/2)}
  \fontsize{6.8pt}{8.2pt}\selectfont
  \begin{block}{Proof — Partial Fractions and Telescoping Sums}
    \fontsize{6.8pt}{8.2pt}\selectfont
    By the normalization axiom, summing across the entire support $\mathcal{S}_N = \{1, \dots, N\}$ must equal one:
    $$\sum_{x=1}^{N} f_N(x) = c_N \sum_{x=1}^{N} \frac{1}{x(x+1)} = 1$$
    We decompose the rational kernel into partial fractions: $\displaystyle \frac{1}{x(x+1)} = \frac{1}{x} - \frac{1}{x+1}$.
  \end{block}
  \vspace{-0.15cm}
  \pause
  \begin{alertblock}{Summation Expansion and Solving for $c_N$}
    \fontsize{6.8pt}{8.2pt}\selectfont
    Expanding the summation explicitly reveals a telescoping cancellation of intermediate terms:
    $$\sum_{x=1}^{N} \left(\frac{1}{x} - \frac{1}{x+1}\right) = \left(1 - \frac{1}{2}\right) + \left(\frac{1}{2} - \frac{1}{3}\right) + \dots + \left(\frac{1}{N} - \frac{1}{N+1}\right) = 1 - \frac{1}{N+1} = \frac{N}{N+1}$$
    Substituting this closed form into our normalization identity:
    $$c_N \left(\frac{N}{N+1}\right) = 1 \implies \mathbf{c_N = \frac{N+1}{N} \quad \blacksquare}$$
  \end{alertblock}
\end{frame}

% Slide 13A: Classroom Exercise — Tier 4: Challenging (Statement)
\begin{frame}{Classroom Exercise — Tier 4: Challenging (1/2)}
  \begin{block}{Problem 3.1.10 — First Failure Time in Parallel Systems}
    \footnotesize
    An engineering system composed of two independent parallel subsystems operates in discrete time cycles $n \in \{1, 2, 3, \dots\}$. In each cycle, subsystem $A$ fails with probability $p_A \in (0,1)$ independently of previous cycles, while subsystem $B$ fails with probability $p_B \in (0,1)$. Let $N_A$ be the first failure cycle of $A$, and $N_B$ the first failure cycle of $B$. We define $T = \min(N_A, N_B)$ as the time until the system experiences its first failure in either branch.
    \begin{enumerate}\itemsep=0.12cm
      \item Derive analytically the exact PMF $f_T(n) = P(T = n)$ for all $n \in \{1, 2, 3, \dots\}$.
      \item Prove that $T$ follows a geometric distribution and explicitly identify its combined failure parameter $p_T$.
    \end{enumerate}
  \end{block}
\end{frame}

% Slide 13B: Classroom Exercise — Tier 4: Challenging (Solution)
\begin{frame}{Classroom Exercise — Tier 4: Challenging (2/2)}
  \fontsize{6.5pt}{7.8pt}\selectfont
  \begin{block}{1. Survival Analysis and Independence of Branches}
    \fontsize{6.5pt}{7.8pt}\selectfont
    The survival probability of each subsystem beyond cycle $n$ requires no failure in $n$ consecutive cycles: $P(N_A > n) = (1-p_A)^n$ and $P(N_B > n) = (1-p_B)^n$. By independence across subsystems:
    $$P(T > n) = P(\min(N_A, N_B) > n) = P(N_A > n \cap N_B > n) = (1-p_A)^n (1-p_B)^n = [(1-p_A)(1-p_B)]^n$$
    To extract the point probability at cycle $n$, we difference consecutive survival probabilities:
    $$f_T(n) = P(T = n) = P(T > n-1) - P(T > n) = [(1-p_A)(1-p_B)]^{n-1} [1 - (1-p_A)(1-p_B)]$$
  \end{block}
  \vspace{-0.15cm}
  \pause
  \begin{alertblock}{2. Identification of the Equivalent Geometric Distribution}
    \fontsize{6.5pt}{7.8pt}\selectfont
    Define the combined probability that at least one subsystem fails in any single cycle:
    $$p_T = 1 - (1-p_A)(1-p_B) = p_A + p_B - p_A p_B$$
    Substituting $p_T$ into our derived mass expression yields:
    $$f_T(n) = (1 - p_T)^{n-1} p_T, \quad n \in \{1, 2, 3, \dots\}$$
    This rigorously proves that $\mathbf{T \sim \text{Geom}(p_T = p_A + p_B - p_A p_B) \quad \blacksquare}$.
  \end{alertblock}
\end{frame}

% Slide 14: Synthesis and Conclusion
\begin{frame}{Synthesis and Conclusion — PMF and Discrete Support}
  \fontsize{7.5pt}{9pt}\selectfont
  \begin{block}{Core Pillars of Section 03.01}
    \fontsize{7.5pt}{9pt}\selectfont
    \begin{itemize}\itemsep=0.12cm
      \item \textbf{Measurable Projection:} A discrete random variable maps outcomes from an arbitrary sample space $\Omega$ onto a finite or countable subset of the real line $\mathbb{R}$.
      \item \textbf{Operational Support:} The support $\mathcal{S}_X$ exclusively contains points with strictly positive mass. Outside $\mathcal{S}_X$, point probability is identically zero.
      \item \textbf{Kolmogorov Axioms:} Every legitimate PMF must unconditionally satisfy non-negativity ($p_X(x) \ge 0$) and unitary normalization over the support ($\sum p_X(x) = 1$).
      \item \textbf{Preimage Transformations:} If $Y=g(X)$, probability mass at $y \in \mathcal{S}_Y$ is computed by summing across all original roots in $g^{-1}(y)$.
    \end{itemize}
  \end{block}
\end{frame}

% Slide 15: Modular Transition to the Next Topic
\begin{frame}{Modular Transition to the Next Topic}
  \fontsize{7.5pt}{9pt}\selectfont
  \begin{block}{From Point Mass to Cumulative Probability}
    \fontsize{7.5pt}{9pt}\selectfont
    We have established the formal foundation for discrete random variables via the Probability Mass Function $p_X(x)$, which answers equality queries ($X=x$).
    \begin{itemize}\itemsep=0.12cm
      \item In engineering design, network traffic management, and reliability analysis, however, critical queries are cumulative threshold events: \emph{``What is the probability of observing at most $k$ failures?''}
      \item To answer such queries compactly, we must accumulate point mass over integer intervals across the support.
    \end{itemize}
  \end{block}
  \vspace{-0.15cm}
  \pause
  \begin{alertblock}{Next Stop: Section 03.02}
    \fontsize{7.5pt}{9pt}\selectfont
    In the next section (\textbf{Section 03.02: Cumulative Distribution Function — Discrete CDF}), we construct the step function $F_X(x) = P(X \le x)$, analyze its jump discontinuities ($\Delta F$), and master the analytical duality between discrete differencing and probability accumulation.
  \end{alertblock}
\end{frame}

\end{document}
